% LUA101 midterm exam: modules 1-7, 75 minutes, 75 points.
%
% One source gives both versions. From this folder:
%   latexmk -lualatex midterm.tex
%     the exam, without the solutions: midterm.pdf
%   latexmk -lualatex -jobname=midterm-key -usepretex='\PassOptionsToClass{answers}{exam}' midterm.tex
%     the key, with every solution and its grading notes: midterm-key.pdf
% The latexmkrc in this folder adds ../../latex to TEXINPUTS, for the shared preamble. Both PDFs
% stay out of Git. This source holds the key, so it lives in solutions/, never in exams/.
\documentclass[11pt,addpoints]{exam}
\usepackage{preamble}
\usepackage[margin=1in]{geometry}
\usepackage{booktabs}

\pagestyle{headandfoot}
\firstpageheader{LUA101: Lua from Front to Back}{}{Midterm Exam}
\runningheader{LUA101}{Midterm Exam}{Page \thepage\ of \numpages}
\firstpagefooter{}{Page \thepage\ of \numpages}{}
\runningfooter{}{}{}

\renewcommand{\solutiontitle}{\noindent\textbf{Solution:}\par\noindent}
\unframedsolutions
\SolutionEmphasis{\color{black!75}}

\lstset{style=lua}

% Each question opens with its number and its total points, parts included.
\qformat{\textbf{Question \thequestion}\quad(\totalpoints\ points)\hfill}

\begin{document}

\begin{center}
  {\Large\bfseries LUA101: Lua from Front to Back}\\[0.5ex]
  {\large Midterm Exam: Modules 1--7}
\end{center}

\vspace{1em}
\noindent Name: \makebox[0.6\textwidth]{\hrulefill}

\vspace{1em}
\noindent\textbf{Instructions.}
\begin{itemize}
  \item You have 75 minutes. The exam has \numquestions\ questions worth \numpoints\ points. Each
    question gives its points, and the learning objectives it assesses; where a question has
    parts, each part gives its share.
  \item Use no computer, phone, or other device. Write each answer in the space provided.
  \item All code is Lua 5.4. When a question asks what a line prints, write exactly what it
    prints; where \code|print| separates values with a tab, a space is enough.
  \item If an expression or a line raises an error, write ``error'' and say why.
\end{itemize}

\begin{center}
  \gradetable[h][pages]
\end{center}

\begin{questions}

% ---------------------------------------------------------------------------------------------
\question[10] \textit{(LO-1.2)} Give the value of each expression, and its type, written as
\code|integer| or \code|float| for a number. If it raises an error, write ``error'' and the
reason.

\begin{center}
\renewcommand{\arraystretch}{1.4}
\begin{tabular}{@{}l l p{1.6in} p{1.6in}@{}}
  \toprule
  & Expression & Value & Type \\
  \midrule
  (a) & \code|7 // 2| & \ifprintanswers 3 \else \fi & \ifprintanswers integer \else \fi \\
  (b) & \code|7 / 2| & \ifprintanswers 3.5 \fi & \ifprintanswers float \fi \\
  (c) & \code|-7 // 2| & \ifprintanswers $-4$ \fi & \ifprintanswers integer \fi \\
  (d) & \code|7 % -3| & \ifprintanswers $-2$ \fi & \ifprintanswers integer \fi \\
  (e) & \code|2^2| & \ifprintanswers 4.0 \fi & \ifprintanswers float \fi \\
  (f) & \code|"10" + 1| & \ifprintanswers 11 \fi & \ifprintanswers integer \fi \\
  (g) & \code|10 .. ""| & \ifprintanswers \code|"10"| \fi & \ifprintanswers string \fi \\
  (h) & \code|"10" == 10| & \ifprintanswers \code|false| \fi & \ifprintanswers boolean \fi \\
  (i) & \code|1 // 0| & \ifprintanswers error \fi & \ifprintanswers (divide by zero) \fi \\
  (j) & \code|7.0 // 2| & \ifprintanswers 3.0 \fi & \ifprintanswers float \fi \\
  \bottomrule
\end{tabular}
\end{center}

\begin{solution}
One point for each row, for the value and the type together. \code|//| floors toward negative
infinity, so \code|-7 // 2| is $-4$, and \code|%| takes the sign of the divisor, so \code|7 % -3| is
$7 - (-3)(-3) = -2$. \code|/| and \code|^| always give floats. Arithmetic converts the string in
(f), but \code|==| never converts, so (h) is \code|false|. Integer floor division by zero raises
an error, while float floor division would give infinity. Common mistakes: 3 for (j), which is a
float because one operand is, and \code|-3| for (c).
\end{solution}

% ---------------------------------------------------------------------------------------------
\question[5] \textit{(LO-1.4)} Give the value of each expression.

\begin{parts}
  \part \code|nil or "x"| \hfill \ifprintanswers \code|"x"| \else \makebox[1.5in]{\hrulefill} \fi
  \part \code|false and nil| \hfill \ifprintanswers \code|false| \else \makebox[1.5in]{\hrulefill} \fi
  \part \code|0 and "zero"| \hfill \ifprintanswers \code|"zero"| \else \makebox[1.5in]{\hrulefill} \fi
  \part \code|"" or 1| \hfill \ifprintanswers \code|""| \else \makebox[1.5in]{\hrulefill} \fi
  \part \code|not 0| \hfill \ifprintanswers \code|false| \else \makebox[1.5in]{\hrulefill} \fi
\end{parts}

\begin{solution}
One point each. Only \code|nil| and \code|false| are false, so 0 and the empty string are true.
\code|and| returns its first operand if it is false, and otherwise its second; \code|or| returns
its first operand if it is true, and otherwise its second. The common mistakes are \code|false|
for (c) and 1 for (d), from languages where 0 and the empty string are false.
\end{solution}

% ---------------------------------------------------------------------------------------------
\question \textit{(LO-1.1)} Running \code|lua inventory.lua| stops with this message:

\begin{lstlisting}[language={},basicstyle=\ttfamily\small]
lua: inventory.lua:27: attempt to index a nil value (local 'row')
\end{lstlisting}

\begin{parts}
  \part[1] Which file and which line raised the error?
  \begin{solutionorbox}[0.5in]
    Line 27 of \code|inventory.lua|.
  \end{solutionorbox}
  \part[1] What was the program doing on that line when it failed, and with which variable?
  \begin{solutionorbox}[0.6in]
    It was reading a field or an item of (indexing) the local variable \code|row|, whose value
    was \code|nil|.
  \end{solutionorbox}
  \part[1] Give one way \code|row| could have been \code|nil| there.
  \begin{solutionorbox}[0.6in]
    Any one cause earns the point, for example: \code|row| came from \code|rows[i]| with an index
    past the end of the list; it came from a misspelled key; or a function that should have
    returned it returned nothing.
  \end{solutionorbox}
\end{parts}

% ---------------------------------------------------------------------------------------------
\question \textit{(LO-1.3)} This function is meant to return the mean of a list.

\begin{lstlisting}
local function average(list)
  local total = 0
  for _, v in ipairs(list) do
    totl = total + v
  end
  return total / #list
end
print(average({ 2, 4, 6 }))
\end{lstlisting}

\begin{parts}
  \part[1] What does the last line print?
  \begin{solutionorbox}[0.5in]
    \code|0.0|
  \end{solutionorbox}
  \part[2] Line 4 has a typing mistake, yet Lua raises no error. Explain why, and what the line
    does instead.
  \begin{solutionorbox}[1in]
    Assigning to a name that is not a declared local creates or changes a global. So on every
    pass, line 4 sets a new global named \code|totl|, and never changes the local \code|total|,
    which stays 0. Dividing 0 by 3 with \code|/| gives the float \code|0.0|.
    One point for ``it creates a global'', and one for explaining why \code|total| stays 0.
  \end{solutionorbox}
  \part[1] Correct the line.
  \begin{solutionorbox}[0.5in]
    \code|total = total + v|
  \end{solutionorbox}
\end{parts}

% ---------------------------------------------------------------------------------------------
\question[5] \textit{(LO-2.1)} For (a) to (d), give the number of times the loop's body runs.

\begin{parts}
  \part \code|for i = 1, 10, 3 do ... end| \hfill \ifprintanswers 4 (1, 4, 7, 10) \else \makebox[1in]{\hrulefill} \fi
  \part \code|for i = 10, 1 do ... end| \hfill \ifprintanswers 0 \else \makebox[1in]{\hrulefill} \fi
  \part \code|for i = 5, 1, -2 do ... end| \hfill \ifprintanswers 3 (5, 3, 1) \else \makebox[1in]{\hrulefill} \fi
  \part \code|for x = 0, 1, 0.25 do ... end| \hfill \ifprintanswers 5 \else \makebox[1in]{\hrulefill} \fi
  \part What does this code print? Explain in one sentence.
\begin{lstlisting}
local i = 100
for i = 1, 3 do end
print(i)
\end{lstlisting}
  \begin{solutionorbox}[0.7in]
    \code|100|. The loop's control variable is a new local of the loop, so the outer \code|i| never
    changes.
  \end{solutionorbox}
\end{parts}

\begin{solution}
One point each. In (b) the default step is 1, so a start above the limit runs no pass. In (d),
0.25 is exact in binary, so the last pass has $x = 1.0$ exactly.
\end{solution}

% ---------------------------------------------------------------------------------------------
\question[6] \textit{(LO-2.2)} With \code|local function f() return 1, 2, 3 end|, answer each
part.

\begin{parts}
  \part \code|#{ f(), f() }| \hfill \ifprintanswers 4 \else \makebox[1.5in]{\hrulefill} \fi
  \part \code|#{ f(), 10 }| \hfill \ifprintanswers 2 \else \makebox[1.5in]{\hrulefill} \fi
  \part \code|select("#", (f()))| \hfill \ifprintanswers 1 \else \makebox[1.5in]{\hrulefill} \fi
  \part \code|select("#", f(), nil)| \hfill \ifprintanswers 2 \else \makebox[1.5in]{\hrulefill} \fi
  \part What \code|print(f(), "x")| prints \hfill \ifprintanswers \code|1 x| \else \makebox[1.5in]{\hrulefill} \fi
  \part The value of \code|d| after \code|local a, b, c, d = f()| \hfill \ifprintanswers \code|nil| \else \makebox[1.5in]{\hrulefill} \fi
\end{parts}

\begin{solution}
One point each. A call gives all its results only as the last item in a list of expressions;
anywhere else, Lua adjusts it to one result, and parentheses also adjust it to one. So (a) is
$1 + 3$ items, and (b), (d), and (e) each use only the first result of the first call. In (f),
the missing fourth value is \code|nil|.
\end{solution}

% ---------------------------------------------------------------------------------------------
\question[4] \textit{(LO-2.3)} For each \code|return| statement in a function, say whether it
is a proper tail call, and give the reason in a few words.

\begin{parts}
  \part \code|return g(x)|
  \begin{solutionorbox}[0.4in] A proper tail call: \code|return| followed by one call. \end{solutionorbox}
  \part \code|return g(x) + 1|
  \begin{solutionorbox}[0.4in] Not a tail call: the caller adds 1 after \code|g| returns. \end{solutionorbox}
  \part \code|return (g(x))|
  \begin{solutionorbox}[0.4in] Not a tail call: the parentheses adjust the results to one after \code|g| returns. \end{solutionorbox}
  \part \code|return x > 0 and g(x)|
  \begin{solutionorbox}[0.4in] Not a tail call: the call is an operand of \code|and|, not the whole expression. \end{solutionorbox}
\end{parts}

\begin{solution}
One point each, for the classification with a correct reason; a classification without a reason
earns half a point.
\end{solution}

% ---------------------------------------------------------------------------------------------
\question \textit{(LO-3.1, LO-3.2)} Consider this code.

\begin{lstlisting}
local t = { "a", "b", "c" }
t[#t + 1] = "d"
t[2] = nil
\end{lstlisting}

\begin{parts}
  \part[2] List every border of \code|t|: every value that \code|#t| is allowed to return.
  \begin{solutionorbox}[0.5in]
    1 and 4. Index 1 is present and index 2 absent; index 4 is present and index 5 absent.
  \end{solutionorbox}
  \part[1] What does \code|for _, v in ipairs(t) do io.write(v) end| print?
  \begin{solutionorbox}[0.4in]
    \code|a|, because \code|ipairs| stops at the first \code|nil|.
  \end{solutionorbox}
  \part[1] What is \code|#{ n = 3, "x", "y" }|?
  \begin{solutionorbox}[0.4in]
    2: the field \code|n| does not count.
  \end{solutionorbox}
  \part[2] Why should a program not rely on \code|#t| for this \code|t|, and what should it do
    instead?
  \begin{solutionorbox}[0.9in]
    The table has a hole, so it is not a sequence, and \code|#t| may return either border
    (Lua 5.4 gives 4 here, but promises nothing). Keep the count in a field such as \code|n|, or
    never store \code|nil| in a list, for example by removing items with \code|table.remove|.
    One point for the reason, one for a remedy.
  \end{solutionorbox}
\end{parts}

% ---------------------------------------------------------------------------------------------
\question[5] \textit{(LO-3.3)} What does the last line print? One point for each value.

\begin{lstlisting}
local a = { 1, 2, 3 }
local b = a
local c = { 1, 2, 3 }
b[1] = 10
local function grow(list) list[#list + 1] = 0 end
local function reset(list) list = {} end
grow(c)
reset(a)
print(a[1], a == b, a == c, #c, #a)
\end{lstlisting}

\begin{solutionorbox}[0.6in]
  \code|10  true  false  4  3|. \code|b| is the same table as \code|a|, so \code|b[1] = 10|
  changes \code|a[1]|, and \code|a == b|. \code|c| is a different table with the same contents, so
  \code|a == c| is \code|false|. \code|grow| changes the table it was passed; \code|reset| only
  points its parameter at a new table, so \code|a| keeps its three items.
\end{solutionorbox}

% ---------------------------------------------------------------------------------------------
\question \textit{(LO-4.2)} Patterns.

\begin{parts}
  \part[1] What does \code|("2026-11-02"):match("(%d+)%-(%d+)")| return?
  \begin{solutionorbox}[0.4in] \code|"2026"| and \code|"11"|: two captures, both strings. \end{solutionorbox}
  \part[1] What does \code|("a.b.c"):gsub(".", "-")| return?
  \begin{solutionorbox}[0.4in] \code|"-----"| and 5: the unescaped dot matches every character. \end{solutionorbox}
  \part[1] What does \code|("<b>x</b>"):match("<(.-)>")| return?
  \begin{solutionorbox}[0.4in] \code|"b"|: the lazy \code|-| stops at the first \code|>|. \end{solutionorbox}
  \part[1] What does \code|("f(a(b)) g(c)"):match("%b()")| return?
  \begin{solutionorbox}[0.4in] \code|"(a(b))"|: the first balanced pair of parentheses, nesting included. \end{solutionorbox}
  \part[4] Write a pattern that matches a whole line of the form \code|name = value| and
    captures the name and the value. The name is one or more letters, digits, or underscores.
    Spaces may appear at the start of the line, around the \code|=|, and at the end, and the
    value must not include the spaces at the end.
  \begin{solutionorbox}[0.9in]
    \code|^%s*([%w_]+)%s*=%s*(.-)%s*$|

    One point each for: anchoring both ends with \code|^| and \code|$|; the name as
    \code|([%w_]+)|; allowing spaces with \code|%s*| in all four places; and a lazy \code|(.-)|
    for the value, so that the final \code|%s*| takes the trailing spaces. A greedy \code|(.*)|
    keeps the trailing spaces and loses the last point.
  \end{solutionorbox}
\end{parts}

% ---------------------------------------------------------------------------------------------
\question \textit{(LO-4.3)} A Python program accepts a word only if the whole word matches the
regular expression \code!(cat|dog)s?!: that is, only the words \code|cat|, \code|cats|,
\code|dog|, and \code|dogs|.

\begin{parts}
  \part[1] Explain why no single Lua pattern can express this test.
  \begin{solutionorbox}[0.6in]
    Lua patterns have no alternation (\verb!|!), and a quantifier can apply to only one
    character class, not to a group of characters.
  \end{solutionorbox}
  \part[3] Write a Lua function \code|is_pet(word)| that returns \code|true| for exactly those
    four words, and \code|false| for every other string.
  \begin{solutionorbox}[1.4in]
\begin{lstlisting}
local pets = { cat = true, dog = true }

local function is_pet(word)
  return pets[word:match("^(.-)s?$")] == true
end
\end{lstlisting}
    The pattern removes at most one final \code|s|, and the set tests the rest. Two points for a
    correct approach, matching the general shape and then testing the varying part, and one for
    the details: \code|cats| is accepted, \code|catss| and the empty string are rejected, and the
    result is a boolean. Any correct code earns full marks, such as comparing \code|word| with
    each of the four words.
  \end{solutionorbox}
\end{parts}

% ---------------------------------------------------------------------------------------------
\question \textit{(LO-5.1)} Consider this code.

\begin{lstlisting}
local function make()
  local n = 0
  return function() n = n + 1; return n end, function() return n end
end
local inc, get = make()
local inc2 = make()
inc(); inc(); inc2()
print(get())                    -- (a)

local fs = {}
for i = 1, 3 do fs[i] = function() return i end end
print(fs[1]() + fs[3]())        -- (b)

local gs, j = {}, 1
while j <= 3 do gs[j] = function() return j end; j = j + 1 end
print(gs[1]() + gs[3]())        -- (c)
\end{lstlisting}

\begin{parts}
  \part[1] What does line (a) print?
  \begin{solutionorbox}[0.4in] \code|2|: \code|inc| and \code|get| share one \code|n|; \code|inc2| has its own. \end{solutionorbox}
  \part[1] What does line (b) print?
  \begin{solutionorbox}[0.4in] \code|4|: $1 + 3$. \end{solutionorbox}
  \part[1] What does line (c) print?
  \begin{solutionorbox}[0.4in] \code|8|: $4 + 4$. \end{solutionorbox}
  \part[2] Explain why (b) and (c) differ.
  \begin{solutionorbox}[0.9in]
    A closure captures a variable, not its value. Each pass of a numeric \code|for| has a new
    local \code|i|, so each closure has its own; the \code|while| loop's \code|j| is one variable,
    shared by all three closures, which is 4 when they run. One point for ``captures the
    variable'', one for the difference between the loops.
  \end{solutionorbox}
\end{parts}

% ---------------------------------------------------------------------------------------------
\question \textit{(LO-6.1)} Consider this function.

\begin{lstlisting}
local function parse_age(text)
  local n = tonumber(text)
  if n == nil then return nil, "not a number: " .. text end
  if n < 0 then error("age must not be negative", 2) end
  return n
end
\end{lstlisting}

\begin{parts}
  \part[1] What does \code|print(parse_age("ten"))| print?
  \begin{solutionorbox}[0.4in] \code|nil  not a number: ten| \end{solutionorbox}
  \part[2] What does \code|print(pcall(parse_age, "-3"))| print? Explain why the message has no
    file name or line number.
  \begin{solutionorbox}[0.8in]
    \code|false  age must not be negative|. Level 2 names the function that called
    \code|parse_age|, and here that is \code|pcall|, a function written in C, which has no line to
    name. One point for the output, one for the reason.
  \end{solutionorbox}
  \part[2] The text comes from a form that a user fills in. Is raising an error for a negative
    age the right choice? Argue from the rule for choosing between an error and a \code|nil|
    result.
  \begin{solutionorbox}[1in]
    No. Raise an error for a mistake in the calling code, and return \code|nil| and a message for
    an expected failure the caller should handle. A user who types a negative age is an expected
    failure, like one who types ``ten'', so both should return \code|nil| and a message. Full
    marks need the rule and its application; the opposite answer earns a point only if it argues
    from the rule, for example that the form has already rejected negative numbers.
  \end{solutionorbox}
\end{parts}

% ---------------------------------------------------------------------------------------------
\question \textit{(LO-7.1, LO-7.2)} Consider this code.

\begin{lstlisting}
local Point = {}
Point.__index = { describe = function(p) return p.x .. "," .. p.y end }
Point.__add = function(a, b)
  return setmetatable({ x = a.x + b.x, y = a.y + b.y }, Point)
end
Point.__eq = function(a, b) return a.x == b.x and a.y == b.y end
local function point(x, y) return setmetatable({ x = x, y = y }, Point) end

local p, q = point(1, 2), point(3, 4)
print((p + q):describe())           -- (a)
print(p == point(1, 2), p == 1)     -- (b)
print(rawget(p, "describe"))        -- (c)
\end{lstlisting}

\begin{parts}
  \part[1] What does line (a) print?
  \begin{solutionorbox}[0.4in] \code|4,6| \end{solutionorbox}
  \part[1] What does line (b) print?
  \begin{solutionorbox}[0.5in]
    \code|true  false|. \code|__eq| compares the coordinates; it runs only when both operands are
    tables, so a table is never equal to a number.
  \end{solutionorbox}
  \part[1] What does line (c) print?
  \begin{solutionorbox}[0.5in]
    \code|nil|: \code|describe| is found through \code|__index|, but it is not stored in
    \code|p|, and \code|rawget| skips metamethods.
  \end{solutionorbox}
  \part[2] Write a metatable \code|frozen| such that assigning any new key to a table with
    this metatable raises the error \code|"no new keys"|, naming the line that made the assignment.
  \begin{solutionorbox}[0.9in]
    \code|local frozen = { __newindex = function() error("no new keys", 2) end }|

    One point for \code|__newindex| as a function that raises the error, one for level 2.
  \end{solutionorbox}
\end{parts}

\end{questions}

\end{document}
